\subsection{Engine}
\paragraph{Descrizione classe:}
Classe che implementa l'intelligenza artificiale del motore scacchistico.
Contiene i metodi per la ricerca della mossa migliore usando algoritmi
di game tree search (Alpha-Beta pruning) e per valutare una posizione.

\paragraph{Costruttori}
Sono presenti due costruttori:

\begin{enumerate}
  \item \texttt{Engine()}: Costruttore vuoto che inizializza l'engine con la posizione iniziale standard
        e profondità di ricerca predefinita di 6 ply.
  \item \texttt{Engine(std::string fen)}: Costruttore che inizializza l'engine con una posizione FEN specifica.
\end{enumerate}

Entrambi i costruttori inizializzano la \textbf{transposition table} globale condivisa.

\paragraph{Attributi pubblici:}
\begin{lstlisting}[language=C++]
chess::Board board;         // Posizione corrente
bool isPlayerWhite;         // Colore del giocatore umano
uint64_t depth;             // Profondita' di ricerca (default: 6)
int64_t eval;               // Valutazione corrente
TTEntry* ttTable;           // Puntatore alla transposition table
static uint64_t nodesSearched; // Contatore nodi esplorati
\end{lstlisting}

\paragraph{Metodi principali:}

\subparagraph{Ricerca e valutazione:}
\begin{itemize}
  \item \texttt{void search(uint64\_t depth)}: Metodo principale che avvia la ricerca della mossa migliore
        fino alla profondità specificata. Usa iterative deepening e aggiorna \texttt{board} con la mossa trovata.

  \item \texttt{int64\_t evaluate(const chess::Board\& board)}: Valuta staticamente una posizione,
        considerando materiale, posizione dei pezzi, mobilità, e sicurezza del re. Restituisce un
        valore positivo se il bianco è in vantaggio, negativo se il nero è in vantaggio.

  \item \texttt{bool isMate()}: Verifica se la posizione corrente è scacco matto.
\end{itemize}

\subparagraph{Generazione e ordinamento mosse:}
\begin{itemize}
  \item \texttt{std::vector<chess::Board::Move> generateLegalMoves(const chess::Board\& b)}:
        Genera tutte le mosse legali per la posizione data. Implementa le ottimizzazioni descritte
        in OPTIMIZATIONS.md (pre-calcolo check/pin).

  \item \texttt{std::vector<ScoredMove> sortLegalMoves(...)}: Ordina le mosse secondo euristiche
        per migliorare il pruning:
        \begin{itemize}
          \item MVV-LVA (Most Valuable Victim - Least Valuable Attacker) per catture
          \item Bonus per promozioni
          \item Bonus per mosse che danno scacco
          \item Killer moves heuristic
          \item History heuristic
          \item Bonus per mosse del re quando sotto scacco
        \end{itemize}

  \item \texttt{chess::Board::Move getBestMove(...)}: Seleziona la mossa con miglior valutazione
        da una lista di mosse.
\end{itemize}

\subparagraph{Utilità valutazione:}
\begin{itemize}
  \item \texttt{int64\_t getMaterialDelta(...)}: Calcola la differenza di materiale tra i due giocatori.
  \item \texttt{int64\_t getMaterialDeltaFAST(...)}: Versione ottimizzata del calcolo materiale.
\end{itemize}

\paragraph{Algoritmo di ricerca:}

L'engine usa un algoritmo di ricerca chiamato \textbf{Negamax con Alpha-Beta pruning}, variante
semplificata del Minimax che sfrutta la simmetria del gioco.

\subparagraph{Componenti principali:}
\begin{enumerate}
  \item \textbf{Alpha-Beta Pruning}: Pota rami dell'albero di ricerca che non possono influenzare
        la decisione finale, riducendo drasticamente i nodi esplorati.

  \item \textbf{Transposition Table}: Hash table globale che memorizza valutazioni di posizioni
        già analizzate per evitare ricalcoli. Usa Zobrist hashing per generare chiavi univoche.
        Ogni entry contiene:
        \begin{lstlisting}[language=C++]
struct TTEntry {
    uint64_t key;    // Zobrist hash della posizione
    int64_t depth;   // Profondita' di ricerca
    int64_t score;   // Valutazione
    uint8_t flag;    // EXACT, LOWERBOUND, o UPPERBOUND
};
        \end{lstlisting}

  \item \textbf{Iterative Deepening}: Ricerca progressivamente più profonda (1-ply, 2-ply, ..., N-ply),
        permettendo di avere sempre una mossa disponibile anche se il tempo scade.

  \item \textbf{Move Ordering}: Ordina le mosse prima della ricerca per migliorare il pruning.
        Mosse promettenti vengono esplorate prima, aumentando la probabilità di beta-cutoff.

  \item \textbf{Killer Moves}: Memorizza mosse non-catture che hanno causato beta-cutoff,
        dandogli priorità in posizioni simili. Array \texttt{killerMoves[2][MAX\_PLY]}.

  \item \textbf{History Heuristic}: Tiene traccia di quanto spesso una mossa (from $\rightarrow$ to)
        ha causato beta-cutoff, usando un array \texttt{history[2][64][64]} (colore × from × to).

  \item \textbf{Late Move Pruning}: Pota mosse "silenziose" (non-catture, non-scacchi) in coda
        alla lista quando a bassa profondità, assumendo che mosse ben ordinate siano già state analizzate.
\end{enumerate}

\subparagraph{Metodo searchPosition:}
Il cuore della ricerca è il metodo privato:
\begin{lstlisting}[language=C++]
int64_t searchPosition(chess::Board& b, int64_t depth,
                       int64_t alpha, int64_t beta, int ply);
\end{lstlisting}

Pseudocodice semplificato:
\begin{verbatim}
searchPosition(board, depth, alpha, beta, ply):
    if depth == 0:
        return evaluate(board)

    # Probe transposition table
    if trovato_in_TT and depth_TT >= depth:
        return valore_TT

    mosse = generateLegalMoves(board)
    if mosse.empty():
        if in_scacco: return -MATE + ply
        else: return 0  # stallo

    ordina_mosse(mosse)  # Move ordering

    for mossa in mosse:
        if late_move_pruning(mossa): continue

        board.movePiece(mossa)
        score = -searchPosition(board, depth-1, -beta, -alpha, ply+1)
        board.unmovePiece(mossa)

        if score >= beta:
            aggiorna_killer_e_history(mossa)
            return beta  # Beta cutoff

        if score > alpha:
            alpha = score
            mossa_migliore = mossa

    salva_in_TT(board, depth, alpha)
    return alpha
\end{verbatim}

\paragraph{Funzione di valutazione:}

La funzione \texttt{evaluate()} combina diversi fattori:

\begin{enumerate}
  \item \textbf{Materiale}: Valore intrinseco dei pezzi (pedone=100, cavallo=320, alfiere=330,
        torre=500, regina=900). Usa \texttt{getMaterialDelta()}.

  \item \textbf{Piece-Square Tables}: Bonus/penalty in base alla posizione di ogni pezzo sulla scacchiera.
        Tavole diverse per apertura/mediogioco ed endgame. Definite in \texttt{piecevaluetables.hpp}.

  \item \textbf{Mobilità}: Numero di mosse legali disponibili (più mosse = migliore).

  \item \textbf{King Safety}: Penalità se il re è esposto, bonus se protetto da pedoni.

  \item \textbf{Struttura pedonale}: Penalità per pedoni doppi, isolati; bonus per pedoni passati.

  \item \textbf{Altri fattori}: Alfiere in coppia, torri su colonne aperte, ecc.
\end{enumerate}

Valori bonus/penalty definiti in \texttt{basebonuspenaltyvalues.hpp}.

\paragraph{Parallelizzazione:}

L'engine usa \textbf{OpenMP} per parallelizzare la ricerca al primo livello (root):
\begin{lstlisting}[language=C++]
#pragma omp parallel for
for (move : root_moves) {
    searchPosition(move, depth-1, ...);
}
\end{lstlisting}

Questo sfrutta CPU multi-core per esplorare diverse mosse candidate contemporaneamente.

\paragraph{Statistiche (modalità DEBUG):}

In build DEBUG, l'engine raccoglie statistiche:
\begin{lstlisting}[language=C++]
static uint64_t ttProbes;      // Interrogazioni TT
static uint64_t ttHits;        // Hit totali TT
static uint64_t ttExactHits;   // Hit con valore esatto
static uint64_t ttCutoffHits;  // Hit con cutoff
\end{lstlisting}

\paragraph{Decisioni progettuali:}

\subparagraph{Negamax vs Minimax:}
Negamax semplifica il codice sfruttando il fatto che \texttt{max(a,b) = -min(-a,-b)}.
Un solo metodo di ricerca invece di due separati per MIN e MAX.

\subparagraph{Transposition Table globale:}
La TT è condivisa tra tutte le istanze di Engine per massimizzare il riuso delle valutazioni.
Dimensione fissa: \texttt{TTEntry::TABLE\_SIZE}.

\subparagraph{Iterative Deepening:}
Anche se sembra controintuitivo ricercare più volte, l'ID permette di:
\begin{itemize}
  \item Avere sempre una mossa disponibile (se interrotto)
  \item Migliorare move ordering usando risultati delle iterazioni precedenti
  \item Nel complesso, essere più veloce grazie al miglior pruning
\end{itemize}

\subparagraph{Costante MAX\_PLY:}
Limite massimo di profondità (64 ply) per evitare stack overflow e dimensionare array statici
come \texttt{killerMoves[2][MAX\_PLY]}.

\subparagraph{Struct helper:}
Le struct \texttt{SearchContext}, \texttt{AlphaBeta}, \texttt{TTSaveInfo} raggruppano
parametri correlati per ridurre il numero di argomenti dei metodi.
